% ===== §10 Eudaimonistic Reconstruction =====
\section{Eudaimonistic Reconstruction: Happiness as Emergent Signal of Relational Co-evolution}
\label{sec:reconstruction}

\noindent
The phenomenological description of the preceding section has shown us, from the inside, what relational happiness is like: how it arises from the between, how it has the character of resonance, how it transforms the quality of time, and how it differs in its deepest form from the intensity of early relational encounter. We are now in a position to give this phenomenological description its theoretical articulation: to state, precisely and at the appropriate level of abstraction, what relational happiness \emph{is}---what kind of thing it is, where it is located, what generates it, and what its relationship is to the individual flourishing that the eudaimonist tradition has described.

This section executes the eudaimonistic reconstruction in three movements. The first states the central claim of the GRB account of happiness---happiness as the emergent signal of relational co-evolution---and defends it against obvious objections. The second locates relational happiness within the three registers of the GRB framework (imaginary, symbolic, real), arguing that the deepest form of relational happiness is a phenomenon of the real register. The third addresses the relationship between relational and individual flourishing, arguing that the GRB account does not eliminate individual flourishing but reconceives it as a moment within the larger structure of relational \emph{eudaimonia}.

\subsection{The Central Claim: Happiness as Emergent Signal}
\label{subsec:emergentsignal}

\noindent
The central claim of this paper's eudaimonistic reconstruction can be stated in a single sentence: \emb{relational happiness is the emergent signal of the relational field's co-evolutionary activity}. Each element of this formulation requires unpacking.

\emb{Emergent.} The happiness is emergent in the technical sense: it is a property of the relational field that cannot be reduced to the properties of the individual systems that compose the field, and that arises at the level of the field through the dynamical interactions of those systems. It is not the sum of two individual happinesses; it is not the average; it is not any function of the individual happinesses that could be computed from knowledge of each individual's state alone. It is a new property, arising at a new level of organisation, irreducible to what exists at the level below.

The emergent character of relational happiness explains why it \emph{feels} as though it comes from the between: because it does. The felt quality of exteriority---the sense that the happiness is coming from somewhere that is not quite within oneself---is an accurate perception of an emergent property. The individual is perceiving, from within the relational field, a property of the field that is not a property of themselves individually. This perception is not an illusion; it is the most accurate perception available of what is actually happening.

\emb{Signal.} The happiness is a signal---a piece of information---that the relational field generates and transmits to the individuals who participate in it. What information does it carry? It carries information about the state of the relational field: specifically, it signals that the field is currently engaged in productive co-evolutionary activity---that the coupling is strong, that the shared attractor landscape is being traversed in a generative way, that the relational STDP mechanism is producing structural modifications that deepen and enrich the field. The happiness is, in this sense, the relational field's way of telling its participants that it is doing well---that it is flourishing, in its own relational mode of flourishing.

This signal character of happiness is not an ad hoc addition to the theory but a consequence of its dynamical foundations. In any complex system with multiple levels of organisation, the higher level generates signals that propagate down to the lower level and modify the behaviour of the lower-level components. The signals that a healthy biological organ sends to the organism of which it is a part---the felt signals of comfort, warmth, satiety, pleasure---are the organ's way of communicating to the organism that it is functioning well and that the organism's behaviour is supporting its continued functioning. The relational happiness signal is the relational field's analogue of these biological signals: it communicates to the individual participants that the field is functioning well and that their shared behaviour is supporting its continued co-evolution.

\emb{Of the relational field's co-evolutionary activity.} The happiness signals a specific kind of relational activity: co-evolution. It is not the signal of any pleasant state of the relational field but of the state in which the field is actively engaged in the co-evolutionary process---in the joint processing of contingent events, the dynamic modification of the shared attractor landscape, the accumulation of relational holonomy. A relational field that is stagnant---that has ceased to co-evolve, that is no longer responsive to contingent events, that has hardened into fixed patterns that no longer develop---does not generate the happiness signal even if it is comfortable, stable, and conflict-free. The happiness requires co-evolution, not mere coexistence.

\begin{claim}[Happiness as emergent signal of co-evolution]
Relational happiness is the emergent signal that the relational field generates when it is engaged in productive co-evolutionary activity. It is emergent because it is a property of the field irreducible to the properties of either individual; it is a signal because it carries information about the field's co-evolutionary state to the individuals who participate in it; and it is specifically the signal of co-evolutionary activity, not of any merely comfortable or stable relational state. The felt happiness of the shared flower is the individual's reception of this signal---the way in which the relational field's co-evolutionary activity presents itself to the individuals who constitute it.
\end{claim}

\subsection{Objections and Replies}
\label{subsec:objections}

\noindent
The central claim will face several objections that must be addressed before proceeding.

\emb{Objection 1: The signal account makes happiness too cognitive.} One might object that describing happiness as a signal that carries information makes it sound too cognitive, too much like a piece of data to be processed, and not enough like the warm, immediate, felt quality of actual happiness. Happiness is not the reception of a signal; it is a lived experience with its own irreducible quality.

\emph{Reply}: The signal account does not deny the irreducible felt quality of happiness; it describes the functional role of that quality within the relational system. Just as the felt quality of pain is not reduced by the functional description of pain as a signal that tissue damage has occurred---the phenomenology of pain remains exactly what it is---the felt quality of relational happiness is not reduced by the functional description of happiness as a signal of co-evolutionary activity. The phenomenological description of the previous section and the functional description of the present section are complementary accounts of the same phenomenon at different levels of description. Neither reduces to the other, and both are necessary for a complete account.

\emb{Objection 2: Not all shared happiness involves co-evolution.} One might object that there are forms of shared happiness---two strangers who laugh together at the same joke, or two people who happen to notice the same beautiful thing simultaneously---that are not instances of co-evolutionary dynamics in any interesting sense. The strangers are not a coupled dynamical system; they do not have a shared attractor landscape; their interaction does not produce structural modifications of a relational field. And yet they are happy, together, in a way that seems genuinely relational.

\emph{Reply}: The GRB account does not deny that there are weaker and stronger forms of relational happiness, or that the happiness of strangers who laugh together is a genuine phenomenon. It claims that the \emph{paradigm} of relational happiness---the form that most fully instantiates the concept and that the eudaimonist tradition has consistently failed to account for---is the happiness of an intimate relational field engaged in deep co-evolutionary activity. The happiness of strangers is a real but thin instance of relational happiness: it involves minimal coupling, no shared attractor landscape, and no structural modification. The happiness of intimate partners sharing a flower is a rich instance: it involves strong coupling, a deep shared attractor landscape, and significant structural modification. The theory is not falsified by the existence of thin instances; it is confirmed by the existence of thick ones and by the fact that the thick instances are precisely where the individualist theories most conspicuously fail.

\emb{Objection 3: The happiness might be explained by the individual's pleasure at the other's pleasure.} A more subtle objection holds that the happiness of the shared flower can be fully explained at the individual level by the following mechanism: person A is pleased by the flower, person B notices person A's pleasure and is pleased by it, and person A notices person B's pleasure at A's pleasure and is further pleased, and so on in a cascade of mutual positive affect. On this account, the happiness is genuinely social in origin---it involves the other's response---but it remains individual in structure: each person's happiness is caused by the other's, but each is still an individual happiness caused by an external stimulus.

\emph{Reply}: This objection correctly identifies a mechanism that is genuinely present in the shared flower---the mutual positive affect cascade is real and empirically well-documented. But it fails to account for the emergent character of the happiness: the fact that the cascade produces something that is not the sum of the individual happinesses but a qualitatively different phenomenon. The mutual affect cascade is the mechanism through which the relational field's co-evolutionary activity is initiated; it is not itself the relational happiness but one of the processes through which relational happiness is generated. More fundamentally, the objection presupposes the individualist ontology that the GRB framework rejects: it assumes that the unit of analysis is the individual happiness, and that the relational happiness is composed of individual happinesses in mutual causal relation. The GRB account holds, on the contrary, that the relational happiness is primary---that it is the happiness of the field, and that the individual's experience of it is the field's happiness perceived from a particular constituted perspective.

\subsection{The Three Registers of Relational Happiness}
\label{subsec:threeregisters}

\noindent
The GRB framework's three registers---imaginary, symbolic, and real---provide a further dimension of the reconstruction. Relational happiness can occur in any of the three registers, and the form it takes in each register is distinct. A complete account of relational \emph{eudaimonia} must map these distinctions.

\emb{Imaginary register happiness} is the happiness of identification and idealisation: the happiness of seeing oneself reflected in the beloved's eyes as the person one wishes to be, or of seeing the beloved as the idealised object of one's desire. This is the happiness of the early relation, of falling in love, of the mirror stage of intimate encounter. It is real and intense---the dramatic spikes of early relational co-evolution correspond, in part, to the powerful imaginary investments of idealisation---but it is unstable: the idealised image is always threatened by the reality of the other's irreducible particularity, and the happiness built on imaginary identification is always liable to collapse when the ideal is disappointed. The happiness of the shared flower, when it occurs primarily in the imaginary register, is the happiness of sharing with the idealised beloved---the beauty of the flower is magnified by the imaginary investment in the person with whom it is shared.

\emb{Symbolic register happiness} is the happiness of recognition and acknowledgment: the happiness of having one's existence, one's contributions, and one's significance acknowledged by the other within the framework of the shared symbolic order---the language, the rituals, the social forms through which intimate relations are constituted and sustained. The gift, the anniversary, the declaration of love, the shared joke that belongs to the couple's private language: these are sources of symbolic register happiness, and they are real and important sources of relational flourishing. But they share the limitation of all symbolic goods: they can be forged, performed, or deployed without genuine relational co-evolution, and the happiness they produce is accordingly shallower and more fragile than the happiness of the real register. \pinkref{Paper~XIII}'s analysis of forged trust---the trust produced by the appearance of genuine self-investment without its substance---is directly relevant here: forged symbolic happiness is the appearance of relational co-evolution without its reality \citep{paperxiii}.

\emb{Real register happiness} is the deepest and most stable form of relational happiness: the happiness of an encounter with the irreducible particularity of the other and of the shared moment, an encounter that resists full symbolisation and exceeds the imaginary investments that surround it. The happiness of the shared flower, at its deepest, is a happiness of the real register: it is the happiness of this particular flower, shared with this particular person, at this particular moment in the history of this particular relational field---a happiness that cannot be captured in any general description, that cannot be fully communicated, that leaves a remainder that exceeds all symbolic elaboration. The real register happiness is the happiness of genuine co-presence: the happiness of being with the other as they actually are, not as they appear in the imaginary or as they are constituted in the symbolic.

\begin{claim}[The priority of real register happiness]
Among the three registers of relational happiness, the real register is primary: it is the deepest, most stable, and most irreducible form, the form that most fully instantiates the concept of relational \emph{eudaimonia}. Imaginary and symbolic register happiness are real and valuable, but they are vulnerable to the disillusionment of idealisation and the inflation of symbolic goods respectively. Real register happiness---the happiness of genuine co-presence, of the encounter with the other's irreducible particularity---is not vulnerable to these failures: it cannot be forged, cannot be inflated, and does not depend on the maintenance of an idealised image.
\end{claim}

\subsection{Relational \emph{Eudaimonia} and Individual Flourishing}
\label{subsec:relindiv}

\noindent
The final question the reconstruction must address is the relationship between relational \emph{eudaimonia}---the flourishing of the relational field---and individual flourishing. The GRB account has argued that the primary subject of flourishing is the relational field, not the individual; but this raises an obvious question: what, then, is the status of individual flourishing? Is it eliminated by the GRB account? Is it a mere epiphenomenon of relational flourishing? Or does it retain an independent significance within the framework?

The answer is that individual flourishing is neither eliminated nor reduced to a mere epiphenomenon, but reconceived as a \emph{moment within} the larger structure of relational \emph{eudaimonia}---a necessary moment, without which relational \emph{eudaimonia} would be impossible, but a moment that achieves its fullest significance within the relational whole.

This reconception can be made precise. Individual flourishing---the development of the individual's capacities, the exercise of the individual's virtues, the realisation of the individual's potential---is a necessary condition for the richest relational flourishing: a relational field composed of individuals who are not themselves developing, exercising their capacities, or realising their potential is a field in which the co-evolutionary dynamics are impoverished. Each individual brings to the relational field the resources of their individual development---their particular capacities, their unique perspective, their specific way of noticing and responding to contingent events---and the richness of the relational field's co-evolutionary dynamics depends on the richness of these individual contributions. Individual flourishing feeds relational flourishing.

But the relationship is also reciprocal: relational flourishing feeds individual flourishing. The individual who participates in a richly co-evolving relational field is themselves developed by that participation---their capacities are expanded by the demands of co-evolution, their perspective is enlarged by the encounter with the other's irreducible particularity, their virtues are exercised and deepened by the practices of shared life. The relational field is not merely the arena in which individual flourishing occurs; it is a generative environment that produces individual flourishing as one of its outputs.

The relationship between relational and individual flourishing is therefore neither unidirectional nor competitive but \emb{mutually generative}: individual flourishing enriches the relational field, and the relational field enriches the individuals who participate in it. This mutual generativity is the formal content of what the GRB framework calls the spiral structure of relational development: each turn of the spiral both requires and produces individual development, and each individual development both requires and produces relational development. The spiral is the form of their mutual implication.

\begin{claim}[Relational \emph{eudaimonia} as mutually generative]
Relational \emph{eudaimonia}---the flourishing of the relational field as the primary subject---and individual flourishing are mutually generative rather than competitive. Individual flourishing enriches the relational field by bringing to it the resources of individual development; relational flourishing enriches the individual by generating, through co-evolutionary dynamics, developments that the individual could not have achieved alone. The spiral structure of relational development is the form of this mutual generativity: each turn of the spiral both requires and produces individual development, and each individual development both requires and produces relational development.
\end{claim}

This mutual generativity is the deepest vindication of the relational account of happiness. It shows that the GRB framework does not sacrifice the individual to the relation---it does not ask the individual to dissolve into the relational field for the sake of relational flourishing---but rather shows that the fullest individual flourishing and the fullest relational flourishing are not in competition but in spiral co-implication. The person who develops most fully as an individual is, in the GRB account, the person who participates most richly in a co-evolving relational field; and the relational field that flourishes most fully is the one whose participants are most fully developed as individuals. This is the eudaimonist content of what Buber said philosophically: in the beginning, and in the end, is the relation.

With the eudaimonistic reconstruction complete, we turn to the hermeneutic re-reading of the classical happiness concepts that the reconstruction makes possible. The question is no longer merely whether these concepts are adequate to relational happiness---we have established that they are not---but how they are transformed when the relational field is placed at the centre of the analysis.
